% Check date
% Add number lines

\documentclass[12pt]{article}
\usepackage[margin=1in]{geometry}
\usepackage{amsmath}
\usepackage{amssymb}

\setlength{\parindent}{0pt}
\setlength{\parskip}{8pt}

\title{Week 3: Inequalities}
\author{
	Student: SA\\
	Tutor: Rachel Eglash}
\date{September 21, 2025}

\begin{document}
	
	\maketitle

	\section*{Week 3 Quiz}
        
        \subsection*{3.1: Inequalities with absolute value; solution sets \& graphs}

	    \subsection*{3.2: Literal equations \& rearranging formulas (solve for a variable)}
        
        \subsection*{3.3: Absolute-value equations; modeling simple contexts}
	    
            \newpage

    \section*{Quick Reference: Equations, Definitions, \& Formulas}
            
        \subsection*{S.A.D.D. — Signs • Arithmetic • Distribute • Denominators}
            \begin{itemize}
                \item \textbf{Signs:} When adding or subtracting, keep the sign of the larger absolute value.
                \item \textbf{Arithmetic:} Follow order of operations (PEMDAS).
                \item \textbf{Distribute:} Use the distributive property to eliminate parentheses.
                \item \textbf{Denominators:} Multiply both sides by the least common denominator to eliminate fractions.
            \end{itemize}
        
        \subsection*{Units?}
        
        \subsection*{Test a value}

        \subsection*{Inequality Symbols}
            \begin{itemize}
                \item $<$ means "less than"
                \item $\leq$ means "less than or equal to"
                \item $>$ means "greater than"
                \item $\geq$ means "greater than or equal to"
            \end{itemize}

	    \subsection*{Solving Linear Inequalities}
        	\begin{itemize}
            	\item Solve like an equation, but remember:
                \item \textbf{When multiplying or dividing by a negative number, flip the inequality sign!}
                \item Example: If $-2x < 6$, divide by $-2$: $x > -3$ (sign flipped)
            \end{itemize}
	
	    \subsection*{Compound Inequalities}
        	\textbf{AND:} Both conditions must be true (intersection)\\
            Example: $-2 \leq x < 5$ means $x$ is between $-2$ and $5$\\\\
            \textbf{OR:} At least one condition must be true (union)\\
            Example: $x < -1$ OR $x > 3$ means $x$ is less than $-1$ or greater than $3$

        \subsection*{Absolute Value Inequalities}
            For $|x| < a$ (where $a > 0$): $-a < x < a$\\\\
            For $|x| > a$ (where $a > 0$): $x < -a$ OR $x > a$

        \subsection*{Interval Notation}
            \begin{itemize}
                \item $(a, b)$ means $a < x < b$ (open interval, excludes endpoints)
                \item $[a, b]$ means $a \leq x \leq b$ (closed interval, includes endpoints)
                \item $[a, b)$ means $a \leq x < b$ (half-open interval)
                \item $(-\infty, a)$ means $x < a$
                \item $[a, \infty)$ means $x \geq a$
            \end{itemize}

            \newpage

    \section*{Week 3 Quiz: Inequalities}
	
        \textbf{Instructions:} Show all work. Circle your final answer.
        
        \subsection*{Part A: Linear Inequalities}
            
            \subsubsection*{1) Solve the inequality: $2x - 5 \geq 9$.}

                \begin{enumerate}
                    \item[(a)] Graph the solution set on a number line.
                    \item[(b)] Write the solution in interval notation.
                \end{enumerate}
                
                \vspace{0.5cm}

            \subsubsection*{2) Solve the inequality: $-3x + 4 < 1$.}

                \begin{enumerate}
                    \item[(a)] Graph the solution set on a number line.
                    \item[(b)] Write the solution in interval notation.
                \end{enumerate}
                
                \vspace{0.5cm}
        
        \subsection*{Part B: Compound Inequalities}
            
            \subsubsection*{3) Solve the compound inequality: $-2 \leq x + 1 < 5$.}

                \begin{enumerate}
                    \item[(a)] Graph the solution set on a number line.
                    \item[(b)] Write the solution in interval notation.
                \end{enumerate}
                
                \vspace{0.5cm}

            \subsubsection*{4) Solve the compound inequality: $2x - 5 \geq 9$ OR $x + 3 < -1$.}

                \begin{enumerate}
                    \item[(a)] Graph the solution set on a number line.
                    \item[(b)] Write the solution in interval notation.
                \end{enumerate}
                
                \vspace{0.5cm}

                \newpage
        
        \subsection*{Part C: Absolute Value Inequalities}
            
            \subsubsection*{5) Solve the inequality: $|x - 3| < 5$.}

                \begin{enumerate}
                    \item[(a)] Graph the solution set on a number line.
                    \item[(b)] Write the solution in interval notation.
                \end{enumerate}
                
                \vspace{0.5cm}

            \subsubsection*{6) Solve the inequality: $|2x + 1| \geq 7$.}

                \begin{enumerate}
                    \item[(a)] Graph the solution set on a number line.
                    \item[(b)] Write the solution in interval notation.
                \end{enumerate}
                
                \vspace{0.5cm}
        
        \subsection*{Part D: Word Problems}
        
            \subsubsection*{7) A rideshare charges a \$3 base fee plus \$0.60 per mile.}
            
                You have at most \$15 to spend.\\
                Let $d =$ miles traveled (in miles).\\
                Assume $d \geq 0$.
                
                \begin{enumerate}
                    \item Write and solve an inequality for $d$.
                    \item State the greatest whole number of miles you can travel.
                    \item[(a)] Graph the solution set on a number line.
                    \item[(b)] Write the solution in interval notation.
                \end{enumerate}
                
                \vspace{0.5cm}

            \subsubsection*{8) A resistor is labeled 220 $\Omega$ (ohms).}
            
                It is acceptable if the resistance differs by at most 5 $\Omega$.\\
                Let $R =$ measured resistance (in ohms).
                
                \begin{enumerate}
                    \item Write and solve an inequality for $R$.
                    \item State the acceptable range.
                    \item[(a)] Graph the solution set on a number line.
                    \item[(b)] Write the solution in interval notation.
                \end{enumerate}

                \newpage

        \subsection*{Week 2 Challenge Review Questions}

            \subsubsection*{Solve the equation: $\frac{1}{3}(x - 2) = 0.25x + 7$}
            
                \vspace{0.5cm}

            \subsubsection*{Given the equation: $ax + b = cx + d$, describe the solution set for each case:}

                \textbf{Case 1:} $a = c$ AND $b = d$
                
                \vspace{0.5cm}
                
                \textbf{Case 2:} $a = c$ AND $b \neq d$
                
                \vspace{0.5cm}
                
                \textbf{Case 3:} $a \neq c$ AND $b = d$
                
                \vspace{0.5cm}
                
                \textbf{Case 4:} $a \neq c$ AND $b \neq d$
                
                \vspace{0.5cm}

                \newpage

    \section*{Solution Guide}

        \subsection*{Part A: Linear Inequalities}

            \textbf{1.} Solve the inequality: $2x - 5 \geq 9$.

                \textit{Solution:}\\
                Add 5 to both sides:\\
                $2x - 5 + 5 \geq 9 + 5$\\
                $2x \geq 14$\\
                \\
                Divide by 2:\\
                $x \geq 7$

                \textit{(a) Graph:} Closed circle at 7, shaded to the right $\rightarrow$

                \textit{(b) Interval notation:} \fbox{$[7, \infty)$}

                \vspace{0.5cm}

            \textbf{2.} Solve the inequality: $-3x + 4 < 1$.

                \textit{Solution:}\\
                Subtract 4 from both sides:\\
                $-3x + 4 - 4 < 1 - 4$\\
                $-3x < -3$\\
                \\
                Divide by $-3$ (flip the inequality sign!):\\
                $x > 1$

                \textit{(a) Graph:} Open circle at 1, shaded to the right $\rightarrow$

                \textit{(b) Interval notation:} \fbox{$(1, \infty)$}

                \newpage

        \subsection*{Part B: Compound Inequalities}

            \textbf{3.} Solve the compound inequality: $-2 \leq x + 1 < 5$.

                \textit{Solution:}\\
                Subtract 1 from all three parts:\\
                $-2 - 1 \leq x + 1 - 1 < 5 - 1$\\
                $-3 \leq x < 4$

                \textit{(a) Graph:} Closed circle at $-3$, open circle at 4, shaded between

                \textit{(b) Interval notation:} \fbox{$[-3, 4)$}

                \vspace{0.5cm}

            \textbf{4.} Solve the compound inequality: $2x - 5 \geq 9$ OR $x + 3 < -1$.

                \textit{Solution:}\\
                Solve each inequality separately:\\
                \\
                First inequality: $2x - 5 \geq 9$\\
                $2x \geq 14$\\
                $x \geq 7$\\
                \\
                Second inequality: $x + 3 < -1$\\
                $x < -4$\\
                \\
                Solution: $x < -4$ OR $x \geq 7$

                \textit{(a) Graph:} Open circle at $-4$, shaded left $\leftarrow$; Closed circle at 7, shaded right $\rightarrow$

                \textit{(b) Interval notation:} \fbox{$(-\infty, -4) \cup [7, \infty)$}

                \newpage

        \subsection*{Part C: Absolute Value Inequalities}

            \textbf{5.} Solve the inequality: $|x - 3| < 5$.

                \textit{Solution:}\\
                For $|x - 3| < 5$, we have:\\
                $-5 < x - 3 < 5$\\
                \\
                Add 3 to all parts:\\
                $-5 + 3 < x < 5 + 3$\\
                $-2 < x < 8$

                \textit{(a) Graph:} Open circles at $-2$ and 8, shaded between

                \textit{(b) Interval notation:} \fbox{$(-2, 8)$}

                \vspace{0.5cm}

            \textbf{6.} Solve the inequality: $|2x + 1| \geq 7$.

                \textit{Solution:}\\
                For $|2x + 1| \geq 7$, we have:\\
                $2x + 1 \leq -7$ OR $2x + 1 \geq 7$\\
                \\
                First inequality: $2x + 1 \leq -7$\\
                $2x \leq -8$\\
                $x \leq -4$\\
                \\
                Second inequality: $2x + 1 \geq 7$\\
                $2x \geq 6$\\
                $x \geq 3$\\
                \\
                Solution: $x \leq -4$ OR $x \geq 3$

                \textit{(a) Graph:} Closed circle at $-4$, shaded left $\leftarrow$; Closed circle at 3, shaded right $\rightarrow$

                \textit{(b) Interval notation:} \fbox{$(-\infty, -4] \cup [3, \infty)$}

                \newpage

        \subsection*{Part D: Word Problems}

            \textbf{7.} WP1) A rideshare charges a \$3 base fee plus \$0.60 per mile. You have at most \$15 to spend.

                \textit{Solution:}\\
                Set up the inequality:\\
                Cost = base fee + (cost per mile $\times$ miles)\\
                $3 + 0.60d \leq 15$\\
                \\
                Subtract 3 from both sides:\\
                $0.60d \leq 12$\\
                \\
                Divide by 0.60:\\
                $d \leq 20$\\
                \\
                Since $d \geq 0$, the solution is $0 \leq d \leq 20$\\
                \\
                Greatest whole number of miles: \fbox{20 miles}

                \textit{(a) Graph:} Closed circles at 0 and 20, shaded between

                \textit{(b) Interval notation:} \fbox{$[0, 20]$}

                \vspace{0.5cm}

            \textbf{8.} WP2) A resistor is labeled 220 $\Omega$ (ohms). It is acceptable if the resistance differs by at most 5 $\Omega$.

                \textit{Solution:}\\
                The resistance can differ by at most 5 $\Omega$, so:\\
                $|R - 220| \leq 5$\\
                \\
                This means:\\
                $-5 \leq R - 220 \leq 5$\\
                \\
                Add 220 to all parts:\\
                $215 \leq R \leq 225$\\
                \\
                Acceptable range: \fbox{215 $\Omega$ to 225 $\Omega$}

                \textit{(a) Graph:} Closed circles at 215 and 225, shaded between

                \textit{(b) Interval notation:} \fbox{$[215, 225]$}

                \newpage

        \subsection*{Week 2 Challenge Review Questions}

            \textbf{Solve the equation:} $\frac{1}{3}(x - 2) = 0.25x + 7$

                \textit{Solution:}\\
                Note that $\frac{1}{4} = 0.25$, so:\\
                $0.25(x - 2) = 0.25x + 7$\\
                \\
                Multiply both sides by $3$:\\
                $x - 2 = 0.75x + 21$\\
                \\
                Subtract $0.75x$ from both sides:\\
                $0.25x - 2 = 21$\\
                \\
                Add $2x$ to both sides:\\
                $0.25x = 23$\\
                \\
                Note that $\frac{1}{4} = 0.25$, so:\\
                $\frac{1}{4}x = 23$\\
                \\
                Multiply both sides by $4$:\\
                $x = 92$
                This is a false statement!

                \fbox{Answer: No solution (contradiction)}

                \vspace{0.5cm}

            \textbf{Given the equation:} $ax + b = cx + d$, describe the solution set for each case:

                \textbf{Case 1:} $a = c$ AND $b = d$

                \textit{Solution:}\\
                The equation becomes $ax + b = ax + b$, which is always true.

                \fbox{Answer: All real numbers (identity)}

                \vspace{0.5cm}

                \textbf{Case 2:} $a = c$ AND $b \neq d$

                \textit{Solution:}\\
                The equation becomes $ax + b = ax + d$ where $b \neq d$.\\
                Subtracting $ax$ from both sides gives $b = d$, which contradicts $b \neq d$.

                \fbox{Answer: No solution (contradiction)}

                \vspace{0.5cm}

                \textbf{Case 3:} $a \neq c$ AND $b = d$

                \textit{Solution:}\\
                The equation becomes $ax + b = cx + b$.\\
                Subtracting $b$ from both sides: $ax = cx$\\
                Rearranging: $(a - c)x = 0$\\
                Since $a \neq c$, we have $a - c \neq 0$, so $x = 0$.

                \fbox{Answer: One solution, $x = 0$}

                \vspace{0.5cm}

                \textbf{Case 4:} $a \neq c$ AND $b \neq d$

                \textit{Solution:}\\
                Rearranging: $ax - cx = d - b$\\
                $(a - c)x = d - b$\\
                Since $a \neq c$, we can divide: $x = \frac{d - b}{a - c}$

                \fbox{Answer: One solution, $x = \frac{d - b}{a - c}$}

\end{document}